\documentclass{article}
\usepackage[margin=1in]{geometry}

\begin{document}

\section*{Project Description}

Attached is a frontend project consisting of five buttons (or tabs). Each button corresponds to one of five database tables. When a button is selected, the corresponding table is displayed, showing 10 randomly selected rows.

In addition, there are two additional buttons that allow the user to:
\begin{itemize}
    \item Add a random row to one of the tables.
    \item Remove a random row from one of the tables.
\end{itemize}

The frontend is already configured to communicate with a backend file named \texttt{main.py}. If preferred, you may implement the backend using JavaScript; however, appropriate modifications to the frontend will be required to ensure compatibility.

The Dockerfile in the github repository with the notes should cover your dependencies for this project. Alternatively, you can install the dependencies specified in notes 12 to complete the project.

\section*{Requirements}

You are required to implement a backend in Python consisting of two modules: \texttt{main.py} and \texttt{crud.py}.

\subsection*{\texttt{main.py}}

The \texttt{main.py} file should:

\begin{itemize}
    \item Serve as the entry point of the backend application.
    \item Establish and expose API endpoints to the frontend.
    \item Handle incoming requests for:
    \begin{itemize}
        \item Retrieving table data (limited to 10 randomly selected rows per request).
        \item Triggering insertion of random rows.
        \item Triggering deletion of random rows.
    \end{itemize}
    \item Delegate all database-related operations to functions defined in \texttt{crud.py}.
\end{itemize}

\subsection*{\texttt{crud.py}}

The \texttt{crud.py} module should:

\begin{itemize}
    \item Manage all direct interactions with the database designed in Project 2.
    \item Implement functions for:
    \begin{itemize}
        \item Querying tables and returning random subsets of rows.
        \item Inserting randomly generated rows into tables.
        \item Deleting randomly selected rows from tables.
    \end{itemize}
    \item Serve as the data access layer separating database logic from application logic.
    \item You can extract columns like this:  \begin{verbatim}colnames =[desc[0] for desc in cur.description]\end{verbatim}
    \item Hint: in order for your data to be properly displayed, after extracting columns and rows from table, combine them like this and return:\begin{verbatim}return [dict(zip(colnames, row)) for row in rows]\end{verbatim}
\end{itemize}

\subsection*{Frontend Table Naming Requirement}

The frontend component (\texttt{App.jsx}) contains a constant array that defines the table names used to interact with the backend:

\begin{verbatim}
const TABLES = [...]
\end{verbatim}

\textbf{You are required to modify this array so that the table names exactly match the table names defined in your backend (\texttt{main.py}) and your database schema from Project 2.}

Failure to ensure consistency between:
\begin{itemize}
    \item \texttt{App.jsx} table names
    \item Backend table validation in \texttt{main.py}
    \item Actual database table names
\end{itemize}

will result in failed API requests and incorrect or missing data display.

\subsection*{Important Note}

Ensure that all table identifiers are consistent across the entire stack:
\begin{itemize}
    \item Frontend (\texttt{App.jsx})
    \item Backend (\texttt{main.py})
    \item Database schema (Project 2)
\end{itemize}



\section*{Additional Requirements}

You must also reuse the database migration script developed in Project 2 to ensure proper schema creation and data migration for integration with the frontend system.

\end{document}